\chapter{Code source complet}
\label{annexe:code}

Cette annexe reproduit l'intégralité du code source du projet. Les fichiers
générés automatiquement par Flex et Bison (\code{lex.yy.c}, \code{ilang.tab.c},
\code{ilang.tab.h}) ne sont pas reproduits, étant entièrement déduits de
\code{ilang.l} et \code{ilang.y}.

\section{ast.h}
\lstinputlisting[style=c]{../ast.h}

\section{ast.c}
\lstinputlisting[style=c]{../ast.c}

\section{symtab.h}
\lstinputlisting[style=c]{../symtab.h}

\section{symtab.c}
\lstinputlisting[style=c]{../symtab.c}

\section{ilang.l : analyseur lexical (Flex)}
\lstinputlisting[style=gram]{../ilang.l}

\section{ilang.y : analyseur syntaxique (Bison)}
\lstinputlisting[style=gram]{../ilang.y}

\section{eval.c : évaluateur}
\lstinputlisting[style=c]{../eval.c}

\section{optim.c : optimisations}
\lstinputlisting[style=c]{../optim.c}

\section{codegen.c : back-end assembleur (extension)}
\lstinputlisting[style=c]{../codegen.c}

\section{main.c : point d'entrée}
\lstinputlisting[style=c]{../main.c}

\section{Makefile}
\lstinputlisting{../Makefile}

\section{Programmes de démonstration}

\subsection{fibonacci.ilang}
\lstinputlisting[style=ilang]{../programmes/fibonacci.ilang}

\subsection{pgcd.ilang}
\lstinputlisting[style=ilang]{../programmes/pgcd.ilang}

\subsection{interactif.ilang}
\lstinputlisting[style=ilang]{../programmes/interactif.ilang}

\subsection{flottants.ilang (extension)}
\lstinputlisting[style=ilang]{../programmes/flottants.ilang}

\subsection{boucles.ilang (extension)}
\lstinputlisting[style=ilang]{../programmes/boucles.ilang}

\chapter{Manuel d'installation}
\label{annexe:install}

\section{Prérequis}

Une distribution GNU/Linux (ou WSL sous Windows) disposant de la chaîne C, de
Flex et de Bison. Sur une base Debian/Ubuntu :

\begin{lstlisting}[style=gram]
sudo apt update
sudo apt install build-essential bison flex make
# Optionnel, pour l'extension assembleur 32 bits :
sudo apt install gcc-multilib
\end{lstlisting}

\section{Compilation}

Depuis le répertoire du projet :

\begin{lstlisting}[style=gram]
make          # compile l'interpreteur -> binaire ./ilang
make clean    # supprime les fichiers generes
make test     # compile puis lance tous les programmes de demonstration
make asm      # demonstration du back-end assembleur (extension)
\end{lstlisting}

La compilation ne doit produire aucun avertissement ni conflit.

\chapter{Guide utilisateur}
\label{annexe:guide}

\section{Invocation}

\begin{lstlisting}[style=gram]
./ilang programme.ilang        # execute un fichier source
./ilang                        # mode interactif (Ctrl-D pour executer)
echo 7 | ./ilang interactif.ilang   # alimenter lire() via un tube
./ilang -s programme.ilang     # genere l'assembleur i386 (extension)
./ilang --no-opt programme.ilang    # desactive les optimisations
\end{lstlisting}

\section{Aide-mémoire du langage}

\begin{table}[ht]
\centering
\begin{tabular}{@{}ll@{}}
\toprule
\textbf{Construction} & \textbf{Syntaxe} \\
\midrule
Affectation        & \code{x = 42} \\
Flottant           & \code{r = 2.5} \\
Arithmétique       & \code{+ - * / \%}, \code{-x}, parenthèses \\
Comparaisons       & \code{== != < > <= >=} \\
Condition          & \code{si (c) \{ ... \} sinon \{ ... \}} \\
Boucle tant que    & \code{tant que (c) \{ ... \}} \\
Boucle pour        & \code{pour (i=0; i<n; i=i+1) \{ ... \}} \\
Rupture de boucle  & \code{arrêter} , \code{continuer} \\
Fonction           & \code{fonction f(a,b) \{ ... retourner(e) \}} \\
Affichage / saisie & \code{afficher(e)} , \code{lire()} \\
Commentaires       & \code{// ligne} , \code{/* bloc */} \\
\bottomrule
\end{tabular}
\end{table}

\chapter{Données expérimentales}
\label{annexe:data}

Les mesures ont été obtenues sur AMD Ryzen~5 5500, sous Ubuntu (WSL), meilleur de
trois exécutions. Elles sont reproductibles via les scripts de mesure fournis.

\begin{table}[ht]
\centering
\caption{Données brutes des benchmarks.}
\begin{tabular}{@{}lll@{}}
\toprule
\textbf{Expérience} & \textbf{Paramètre} & \textbf{Mesure} \\
\midrule
Fibonacci (interpréteur) & $n=28$ & 0,43 s \\
Fibonacci (interpréteur) & $n=30$ & 1,13 s \\
Fibonacci (interpréteur) & $n=32$ & 2,97 s \\
Fibonacci (asm i386)     & $n=32$ & 0,01 s \\
Boucle 30~M (avec optim.) & pliage actif   & 4,05 s \\
Boucle 30~M (sans optim.) & \code{--no-opt} & 6,02 s \\
Récursion \code{somme}    & $n=5000$  & ok \\
Récursion \code{somme}    & $n=10000$ & débordement de pile \\
\bottomrule
\end{tabular}
\end{table}

\chapter{Assembleur i386 généré (exemples)}
\label{annexe:asm}

Cette annexe reproduit le code assembleur réellement produit par le back-end
(\code{./ilang -s}) pour deux programmes. Ce code est compilable tel quel avec
\code{gcc -m32 -no-pie} et produit, à l'exécution, exactement les mêmes résultats
que l'interpréteur (respectivement 6 et 120).

\section{PGCD (boucle \code{tant que})}
Source : voir \code{pgcd.ilang} en annexe~\ref{annexe:code}. Le code généré
illustre la traduction d'une boucle en labels et sauts conditionnels.
\lstinputlisting[style=asm]{generated/pgcd.s}

\section{Factorielle (fonction récursive)}
Pour \code{factorielle(5)}, le code généré illustre la récursion : la fonction
\code{il\_factorielle} s'appelle elle-même, chaque appel disposant de son propre
cadre de pile.
\lstinputlisting[style=asm]{generated/factorielle.s}

\chapter{Contenu de l'archive}
\label{annexe:archive}

Le code source est livré dans l'archive
\code{IED-L3-IC-CHAPITRE\_10-ILIAS\_AIT\_BENAISSA-82508880-sources.zip}, dont
l'arborescence est la suivante :

\begin{lstlisting}[style=gram]
ilang/
|-- ilang.l            analyseur lexical (Flex)
|-- ilang.y            grammaire et construction de l'AST (Bison)
|-- ast.h  ast.c       definition et manipulation de l'AST, type Value
|-- symtab.h  symtab.c table des symboles (pile de cadres + hachage)
|-- eval.c             evaluateur recursif
|-- optim.c            optimisations (pliage de constantes, code mort)
|-- codegen.c          back-end assembleur i386 (extension hors-perimetre)
|-- main.c             point d'entree (fichier, interactif, -s, --no-opt)
|-- Makefile           compilation automatisee
|-- README.md          instructions de compilation et d'utilisation
|-- JOURNAL.md         journal de bord (difficultes, choix de conception)
|-- programmes/        programmes de demonstration
|   |-- fibonacci.ilang
|   |-- pgcd.ilang
|   |-- interactif.ilang
|   |-- flottants.ilang     (extension : flottants)
|   |-- boucles.ilang       (extension : arreter / continuer)
|-- tests/             cas de test
    |-- err_variable.ilang   err_divzero.ilang   err_fonction.ilang
    |-- err_arguments.ilang  err_lexical.ilang   err_syntaxe.ilang
    |-- optim.ilang          run_errors.sh
\end{lstlisting}

\paragraph{Remarque sur les fichiers générés.} L'archive ne contient
volontairement \emph{pas} les fichiers produits automatiquement lors de la
compilation, à savoir \code{lex.yy.c} (engendré par Flex à partir de
\code{ilang.l}), ainsi que \code{ilang.tab.c} et \code{ilang.tab.h} (engendrés par
Bison à partir de \code{ilang.y}). Ces fichiers sont entièrement reconstruits par
la commande \code{make} et n'ont donc pas à figurer dans les sources. De même, les
fichiers objets \code{.o} et le binaire \code{ilang} ne sont pas livrés, puisqu'ils
résultent de la compilation.

\paragraph{Récapitulatif des sorties.} Le tableau~\ref{tab:recap} rassemble, pour
mémoire, les sorties attendues des programmes de démonstration, déjà détaillées au
chapitre~\ref{chap:validation}.

\begin{table}[ht]
\centering
\caption{Récapitulatif des sorties des programmes de démonstration.}
\label{tab:recap}
\small
\begin{tabular}{@{}lll@{}}
\toprule
\textbf{Programme} & \textbf{Entrée} & \textbf{Sortie} \\
\midrule
\code{fibonacci.ilang}  & aucune & \code{0 1 1 2 3 5 8 13 21 34 55 89 144 233 377} \\
\code{pgcd.ilang}       & aucune & \code{6} \\
\code{interactif.ilang} & \code{7}  & \code{999} puis \code{49} \\
\code{interactif.ilang} & \code{-5} & \code{999} puis \code{5} \\
\code{flottants.ilang}  & aucune & \code{6.28 / 3.5 / 2 / 2.5 / 19.6349 / 1} \\
\code{boucles.ilang}    & aucune & \code{0 1 2 3 4 / 1 3 5 7 9 / 10 9 8} \\
\bottomrule
\end{tabular}
\end{table}

\chapter{Glossaire}
\label{annexe:glossaire}

\begin{description}[style=nextline,leftmargin=1em]
  \item[AST] \emph{Abstract Syntax Tree}, arbre syntaxique abstrait :
        représentation arborescente du programme, débarrassée des détails de
        syntaxe concrète.
  \item[BNF] \emph{Backus-Naur Form}, notation pour décrire formellement une
        grammaire hors-contexte.
  \item[cdecl] convention d'appel de fonction sur architecture x86 (arguments
        empilés de droite à gauche, nettoyage par l'appelant).
  \item[LALR(1)] \emph{Look-Ahead LR}, classe de grammaires analysables de manière
        ascendante avec un symbole de prévision ; celle qu'accepte Bison.
  \item[Lexème / token] unité lexicale élémentaire reconnue par l'analyseur
        lexical (nombre, identifiant, mot-clé, opérateur).
  \item[PIE] \emph{Position-Independent Executable}, format d'exécutable par
        défaut sur les systèmes modernes, incompatible avec l'adressage absolu
        utilisé par le back-end (d'où l'option \code{-no-pie}).
  \item[Pliage de constantes] optimisation calculant à la compilation les
        sous-expressions entièrement constantes.
  \item[Portée lexicale] règle de visibilité des variables fondée sur la structure
        du programme (et non sur la chaîne dynamique des appels).
  \item[Shift/reduce] type de conflit qu'un analyseur LR peut rencontrer lorsqu'il
        hésite entre décaler un symbole et réduire une règle.
\end{description}
